\section{Related work}

% \vspace{-10pt}
% \begin{table*}[ht]
% \centering
% \caption{
% Comparison of existing work evaluating large language models (LLMs) on hallucination and agentic capabilities. The columns represent the following criteria: whether LLMs are prompted to take actions in the external world (Agentic Setting); whether the context of interaction is in multi-turn and dynamic environments (Multi-turn Interactive Context); whether domains of tasks involve multiple interactive environments (Multiple Agent Env.); whether the setup avoids complex environment configuration (No Env. Setup); whether the context involves realistic risk triggers encountered in the wild (Real-world Risk Trigger); and whether automatic and scalable verification of unfaithful behaviors without manual annotation is supported (Scalable Hallucination Verifier).
% }
% \resizebox{2\columnwidth}{!}{%
% \begin{tabular}{@{}lcccccc@{}}
% \toprule
% & \shortstack{Agentic\\Setting?} 
% & \shortstack{Multi-turn\\Interactive Context?} 
% & \shortstack{Multiple\\Agent Env.?} 
% & \shortstack{No Env.\\Setup?} 
% & \shortstack{Real-world\\Risk Trigger?} 
% & \shortstack{Hallucination\\Verifier?} \\ 
% \midrule
% \textsc{HaluEval-2.0~\cite{li2024dawn}} & \xmark & \xmark & - & - & \xmark & \cmark \\
% \textsc{HaloGEN~\cite{ravichander2025halogen}} & \xmark & \xmark & - & - & \cmark & \cmark \\
% \textsc{EnvDistraction~\cite{ma2024caution}} & \cmark & \xmark & - & - & \cmark & \cmark \\
% \textsc{WebArena~\cite{zhou2023webarena}} & \cmark & \cmark & \xmark & \xmark & \cmark & \xmark \\
% \textsc{TheAgentCompany~\cite{xu2024theagentcompany}} & \cmark & \cmark & \xmark & \xmark & \xmark & \xmark \\
% \textsc{AgentBench~\cite{liu2023agentbench}} & \cmark & \cmark & \cmark & \xmark & \xmark & \xmark \\
% \midrule
% \textsc{\textbf{Ours}} & \cmark & \cmark & \cmark & \cmark & \cmark & \cmark \\
% \bottomrule
% \vspace{-8mm}
% \end{tabular}
% }
% \label{tab:research_comparsion}
% \end{table*}

\vspace{-2mm}
\paragraph{Generalist LLM Agents Benchmark.}
Benchmarks have evaluated large language model (LLM)-based agents on translating high-level instructions into actionable sequences within complex environments. Existing benchmarks such as Mind2Web~\cite{deng2023mind2web} relies primarily on static user interaction traces, thus lacking the ability to capture dynamic behaviors. Later efforts like WebArena, WorkArena, and OSWorld introduced dynamic, self-hosted environments for robust, end-to-end agent evaluations through environmental state analysis~\cite{zhou2023webarena,drouin2024workarena,xie2024osworld}.

Other benchmarks, including TheAgentCompany~\cite{xu2024theagentcompany} and $\tau$-Bench~\cite{yao2024taubenchbenchmarktoolagentuserinteraction}, expanded interactions by integrating simulated human NPCs (non-player characters, supported by LLMs), enabling deeper assessment of agent reasoning, planning, and social behaviors. Similarly, SWE-Bench evaluates agent performance within software engineering contexts, involving interactions with repositories and execution environments~\cite{jimenez2023swe}.

Comprehensive frameworks such as AgentBench~\cite{liu2023agentbench} and BrowserGym~\cite{chezelles2024browsergym} systematically analyze agent capabilities across diverse applications, providing structured evaluations and detailed insights~\cite{bai2024digirl}. Unlike these frameworks that mainly focus on task completion efficacy, \textsc{MIRAGE-Bench} specifically targets the fidelity of agent actions, analyzing the frequency and nature of contextually unfaithful behaviors.

\begin{table}[t]
    \centering
    \resizebox{\columnwidth}{!}{%
    \begin{tabular}{@{}lcccccc@{}}
        \toprule
        & \shortstack{Agentic\\Setting?}
        & \shortstack{Multi-turn\\Interactive Context?}
        & \shortstack{Multiple\\Agent Env.?}
        & \shortstack{No Env.\\Setup?}
        & \shortstack{Real-world\\Risk Trigger?}
        & \shortstack{Hallucination\\Verifier?} \\
        \midrule
        \textsc{HaluEval-2.0~\cite{li2024dawn}} & \xmark & \xmark & - & - & \xmark & \cmark \\
        \textsc{HaloGEN~\cite{ravichander2025halogen}} & \xmark & \xmark & - & - & \cmark & \cmark \\
        \textsc{EnvDistraction~\cite{ma2024caution}} & \cmark & \xmark & - & - & \cmark & \cmark \\
        \textsc{WebArena~\cite{zhou2023webarena}} & \cmark & \cmark & \xmark & \xmark & \cmark & \xmark \\
        \textsc{TheAgentCompany~\cite{xu2024theagentcompany}} & \cmark & \cmark & \xmark & \xmark & \xmark & \xmark \\
        \textsc{AgentBench~\cite{liu2023agentbench}} & \cmark & \cmark & \cmark & \xmark & \xmark & \xmark \\
        \midrule
        \textsc{MIRAGE-Bench~(\textbf{Ours})} & \cmark & \cmark & \cmark & \cmark & \cmark & \cmark \\
        \bottomrule
    \end{tabular}
    }
    \vspace{5pt}
    \caption{\footnotesize
        Comparison of existing work evaluating LLMs on hallucination and agentic capabilities. The columns represent the following criteria: whether LLMs are prompted to take actions in the external world (Agentic Setting); whether the context of interaction is in multi-turn and dynamic environments (Multi-turn Interactive Context); whether domains of tasks involve multiple interactive environments (Multiple Agent Env.); whether the setup avoids complex environment configuration (No Env. Setup); whether the context involves realistic risk triggers encountered in the wild (Real-world Risk Trigger); and whether automatic and scalable verification of unfaithful behaviors without manual annotation is supported (Scalable Hallucination Verifier).
    }
    \label{tab:research_comparsion}
    % \vspace{-10pt}
\end{table}

\paragraph{Hallucination in LLMs.}
The phenomenon of hallucination---where language generation systems produce outputs that are unfaithful to the information available to the model---has long posed a significant challenge in natural language generation (NLG) tasks~\cite{ji2023survey}. In the context of large language models (LLMs), hallucinations have been extensively studied in terms of their taxonomy~\cite{rawte2023troubling,huang2025survey}, benchmarking and evaluation~\cite{li2024dawn,niu2023ragtruth,luo2024halludial}, and the development of detection and mitigation techniques~\cite{farquhar2024detecting,varshney2023stitch, tonmoy2024comprehensive}. HalluLens~\cite{bang2025hallulens} further disentangle LLM hallucination from factuality and defined two types of hallucination: \textbf{Extrinsic hallucination} where generated content is inconsistent with the training data and \textbf{Intrinsic hallucination} where generated content is inconsistent with the contextual input. \textsc{MIRAGE-Bench} further extends the study of hallucinations to agentic settings, an area that remains largely underexplored by existing research. 

\paragraph{Risk of LLM agents.}
Despite growing autonomy, LLM agents remain vulnerable to adversarial manipulations---such as deceptive web elements~\cite{wu2024adversarial,zhang2024attacking} and hidden prompt injections~\cite{xu2024advweb}---that can lead to incorrect or malicious actions. Even benign distractions like pop-ups can trigger unfaithful behaviors, causing agents to deviate from intended tasks or become exploitable~\cite{ma2024caution}. Our work extends the realistic risk consideration beyond single-turn distractions, systematically introducing broader risk settings within dynamic, multi-turn interactions that elicit unfaithful decision-making.

In summary, compared with existing benchmarks in Table~\ref{tab:research_comparsion}, \textsc{MIRAGE-Bench} addresses critical gaps in studying hallucination in LLM agents through authentic interactive context based on realistic environments.